% !PRINTED 367
\weberpage{367}
%<<CONT
und da hier $n = 11$, also ungerade ist, so hat man zu setzen:
\[
x = -5297, \quad y = -781.
\]
Dies sind die absolut kleinsten Werthe von $x$, $y$, die der Gleichung (1) genügen.

Die kleinste positive Lösung erhält man daraus, wenn man $\alpha$ und $\beta$ dazu addirt, also
\[
x = 19038, \quad y = 2807.
\]
Hierauf wird auch die Lösung der allgemeineren Gleichung

\weq{4}{\alpha y - \beta x = \gamma}
in ganzen Zahlen $x$, $y$ zurückgeführt.

Zunächst ist klar, dass, wenn (4) überhaupt lösbar sein soll, jeder gemeinsame Theiler von $\alpha$ und $\beta$ auch Theiler von $\gamma$ sein muss. Ein solcher gemeinsamer Theiler wird dann durch Division weggeschafft. Wir nehmen also auch hier an, dass $\alpha$ und $\beta$ keinen gemeinsamen Theiler haben. Dann folgt ebenso wie oben, dass alle Lösungen von (4) aus einer von ihnen, $x_0$, $y_0$, erhalten werden in der Form
\[
x = x_0 + h \alpha, \quad y = y_0 + h \beta,
\]
worin $h$ eine unbestimmte ganze Zahl ist. Setzt man dann in (4)
\[
x = \gamma \xi, \quad y = \gamma \eta,
\]
und theilt durch $\gamma$, so geht (4) über in
\[
\alpha \eta - \beta \xi = 1,
\]
also in eine Gleichung von der Form (1).

Die Lösung wird zu einer völlig bestimmten, wenn noch eine Bedingung gegeben ist, aus der $h$ bestimmt werden kann, z.\,B. die, dass $y$ zwischen $0$ und $\beta$ liegen soll, mit Einschluss der einen der beiden Grenzen. Wir wollen dies in folgendem Lehrsatz zusammenfassen:

\gesperrt{III. Die unbestimmte Gleichung}
\[
\alpha y - \beta x = \gamma
\]
\gesperrt{hat, wenn $\alpha$, $\beta$, $\gamma$ ganze Zahlen und $\alpha$ und $\beta$ ohne gemeinsamen Theiler sind, immer eine ganzzahlige Lösung $x$, $y$ und nur eine, wenn noch die Bedingung}
\[
0 \leqq y < \beta,
\]
\gesperrt{oder}
\[
0 < y \leqq \beta
\]
\gesperrt{hinzukommt.}
